\documentclass[11pt, a4paper]{article}

% --- UNIVERSAL PREAMBLE BLOCK ---
\usepackage[a4paper, top=2.5cm, bottom=2.5cm, left=2cm, right=2cm]{geometry}
\usepackage{fontspec}

\usepackage[english, bidi=basic, provide=*]{babel}

\babelprovide[import, onchar=ids fonts]{english}

% Set default/Latin font to Sans Serif in the main (rm) slot
\babelfont{rm}{Noto Sans}

% --- ADDITIONAL PACKAGES ---
\usepackage{amssymb}    % For checkboxes \square
\usepackage{booktabs}   % For professional tables
\usepackage{longtable}  % For multi-page tracking table
\usepackage{enumitem}   % For customized lists
\usepackage{xcolor}     % For colored highlights
\usepackage[colorlinks=true, linkcolor=blue]{hyperref} % For internal links

% --- CUSTOM COMMANDS ---
\newcommand{\project}[1]{\vspace{1mm}\noindent\textbf{\textcolor{blue}{Hands-on Project:}} #1}
\newcommand{\usecase}[1]{\vspace{1mm}\noindent\textbf{\textcolor{teal}{Use Case \& Architecture:}} #1}
\newcommand{\capstone}[2]{
    \vspace{3mm}
    \noindent\fbox{
        \parbox{0.97\textwidth}{
            \textbf{\textcolor{red}{CAPSTONE PROJECT - #1:}} #2
        }
    }
    \vspace{3mm}
}
\newcommand{\dayitem}[2]{\item \textbf{#1:} #2}

\title{\Huge \textbf{AWS Master Data Engineer}\\ \Large 6-Month Exhaustive Training \& Certification Program}
\author{Advanced Architecture, Hands-on Labs, and MLOps}
\date{}

\begin{document}

\maketitle

\begin{abstract}
\noindent This 24-week program is designed to transform you from a practitioner to a Master Data Engineer on Amazon Web Services (AWS). It progressively builds from core S3 data lakes to advanced AWS Lake Formation governance, integrating daily theory, weekly service-specific mini-projects, architectural use case analyses, and massive monthly Capstone projects. A tracking timetable is included at the end.
\end{abstract}

\tableofcontents
\newpage

% ==========================================
% MILESTONE 1
% ==========================================
\section{Milestone 1: Core Data Infrastructure \& Storage (Month 1)}
\textit{Objective: Master the foundational storage options, relational databases, and NoSQL.}

\subsection*{Week 1: Amazon S3 \& Storage Foundations}
\begin{itemize}[label=-]
    \dayitem{Monday}{S3 Architecture, Storage Classes (Standard, IA, Glacier Flexible Retrieval/Deep Archive).}
    \dayitem{Tuesday}{Lifecycle Policies, Versioning, Object Lock, and S3 Replication.}
    \dayitem{Wednesday}{AWS DataSync, Storage Gateway, and AWS Snow Family.}
    \dayitem{Thursday}{\project{Deploy an S3 bucket via AWS CLI, configure a 30-day lifecycle transition to Glacier, and set up event notifications to SQS.}}
    \dayitem{Friday}{\usecase{Design a multi-region, compliant disaster recovery storage architecture for a healthcare provider processing 10PB of X-ray images.}}
\end{itemize}

\subsection*{Week 2: Relational Databases (RDS \& Aurora)}
\begin{itemize}[label=-]
    \dayitem{Monday}{Amazon RDS Architecture, Multi-AZ deployments, and Read Replicas.}
    \dayitem{Tuesday}{Amazon Aurora: Cloud-native storage, Global Databases, and Aurora Serverless v2.}
    \dayitem{Wednesday}{Database Migration Service (AWS DMS) and Schema Conversion Tool (SCT).}
    \dayitem{Thursday}{\project{Deploy an Amazon Aurora Serverless v2 cluster, create a cross-region read replica, and simulate a regional failover.}}
    \dayitem{Friday}{\usecase{Architect a zero-downtime migration strategy from a monolithic on-premise Oracle ERP system to Amazon Aurora PostgreSQL.}}
\end{itemize}

\subsection*{Week 3: High-Scale NoSQL (Amazon DynamoDB)}
\begin{itemize}[label=-]
    \dayitem{Monday}{DynamoDB core concepts: Partitions, RCU/WCU (Provisioned vs. On-Demand).}
    \dayitem{Tuesday}{Data Modeling: Primary keys, Sort keys, LSI vs. GSI (Global Secondary Indexes).}
    \dayitem{Wednesday}{DynamoDB Streams, TTL (Time to Live), and Global Tables.}
    \dayitem{Thursday}{\project{Create a DynamoDB table, design a single-table schema for a forum, write a Python Boto3 script to insert items, and configure a GSI for querying.}}
    \dayitem{Friday}{\usecase{Design an ingestion and serving layer for a gaming leaderboard supporting 1 million reads/writes per second globally.}}
\end{itemize}

\subsection*{Week 4: Hybrid Ingestion \& Caching}
\begin{itemize}[label=-]
    \dayitem{Monday}{Amazon ElastiCache (Redis/Memcached) for database acceleration.}
    \dayitem{Tuesday}{DynamoDB Accelerator (DAX).}
    \dayitem{Wednesday}{AWS AppSync (GraphQL) and API Gateway data ingestion.}
    \dayitem{Thursday}{\project{Deploy a Redis cluster via ElastiCache. Write a Python script to cache query results from RDS to reduce database load.}}
    \dayitem{Friday}{\usecase{Architect a high-throughput, low-latency data API for a mobile application that experiences extreme traffic spikes during live sports events.}}
\end{itemize}

\capstone{Milestone 1}{Build an E-Commerce Multi-Tier Data Backend. Deploy Aurora Serverless for user profiles/billing, DynamoDB for active shopping carts and clickstream tracking, S3 for product images, and AWS DMS to continuously replicate Aurora data to S3 for analytics.}

% ==========================================
% MILESTONE 2
% ==========================================
\section{Milestone 2: Data Warehousing \& Analytics (Month 2)}
\textit{Objective: Deep dive into Redshift optimization, Serverless analytics, and Data Lakes.}

\subsection*{Week 5: Amazon Redshift Architecture \& Basics}
\begin{itemize}[label=-]
    \dayitem{Monday}{Redshift MPP architecture, Leader vs. Compute nodes, RA3 nodes.}
    \dayitem{Tuesday}{Redshift Serverless vs. Provisioned pricing and scaling.}
    \dayitem{Wednesday}{Data ingestion: COPY commands, Auto-Copy from S3.}
    \dayitem{Thursday}{\project{Provision a Redshift Serverless endpoint. Load 10 million rows of synthetic JSON data from S3 using the COPY command with IAM roles.}}
    \dayitem{Friday}{\usecase{Estimate the monthly cost and recommend an architecture (RA3 vs. Serverless) for a BI workload that is heavily utilized only during end-of-month reporting.}}
\end{itemize}

\subsection*{Week 6: Redshift Optimization \& Performance}
\begin{itemize}[label=-]
    \dayitem{Monday}{Distribution Styles (EVEN, KEY, ALL, AUTO).}
    \dayitem{Tuesday}{Sort Keys (Compound vs. Interleaved) and Compression Encodings.}
    \dayitem{Wednesday}{Workload Management (WLM), Concurrency Scaling, and Materialized Views.}
    \dayitem{Thursday}{\project{Create three identical tables in Redshift using different Distribution Styles. Benchmark the execution time of a complex JOIN query on each table.}}
    \dayitem{Friday}{\usecase{Re-architect a poorly performing multi-petabyte analytics table that is causing CPU bottlenecks during complex daily dashboard aggregations.}}
\end{itemize}

\subsection*{Week 7: Data Lakehouse \& Amazon Athena}
\begin{itemize}[label=-]
    \dayitem{Monday}{Amazon Athena: Serverless Presto querying directly on S3.}
    \dayitem{Tuesday}{AWS Glue Data Catalog: Crawlers, Databases, and Tables.}
    \dayitem{Wednesday}{Redshift Spectrum: Querying S3 data directly from Redshift.}
    \dayitem{Thursday}{\project{Set up a Glue Crawler to catalog an S3 bucket of raw CSVs. Write an Athena query to aggregate the data, then query the same catalog using Redshift Spectrum.}}
    \dayitem{Friday}{\usecase{Design a cost-effective "Lakehouse" architecture for a media company wanting to store 50PB of logs in S3 while querying them alongside core business data in Redshift.}}
\end{itemize}

\subsection*{Week 8: Data Sharing \& Lake Formation}
\begin{itemize}[label=-]
    \dayitem{Monday}{AWS Lake Formation: Centralized data lake governance.}
    \dayitem{Tuesday}{Row-level and column-level security in Lake Formation.}
    \dayitem{Wednesday}{Redshift Data Sharing across AWS accounts.}
    \dayitem{Thursday}{\project{Configure AWS Lake Formation over an S3 bucket. Create a data permission policy that restricts a "marketing" IAM role from seeing the "PII\_Email" column in Athena.}}
    \dayitem{Friday}{\usecase{Architect a Data Clean Room for an ad agency to merge their CRM data with third-party advertiser data without exposing raw PII to either party.}}
\end{itemize}

\capstone{Milestone 2}{Enterprise Data Lakehouse Build. Ingest 5 years of retail data into S3 using Parquet. Catalog it with AWS Glue. Use Lake Formation to apply row/column level security policies based on department. Create Redshift Materialized Views to serve executive dashboards via Redshift Spectrum.}

% ==========================================
% MILESTONE 3
% ==========================================
\section{Milestone 3: Batch Processing \& ETL Orchestration (Month 3)}
\textit{Objective: Master Spark on EMR, Glue ETL, and Airflow orchestration.}

\subsection*{Week 9: Amazon EMR (Managed Hadoop/Spark)}
\begin{itemize}[label=-]
    \dayitem{Monday}{EMR Architecture, Master/Core/Task nodes, and Spot Instances.}
    \dayitem{Tuesday}{EMR Serverless vs. EMR on EC2 vs. EMR on EKS.}
    \dayitem{Wednesday}{EMRFS: Using S3 as the HDFS storage layer.}
    \dayitem{Thursday}{\project{Deploy an EMR cluster with Spark. Submit a PySpark job that reads data from S3, performs a distributed aggregation, and writes the output back to S3 in Parquet format.}}
    \dayitem{Friday}{\usecase{Design a migration strategy for a 100-node on-premise Cloudera Hadoop cluster, prioritizing cost reduction and decoupling compute from storage in AWS.}}
\end{itemize}

\subsection*{Week 10: Serverless ETL (AWS Glue)}
\begin{itemize}[label=-]
    \dayitem{Monday}{AWS Glue ETL: DynamicFrames vs. Spark DataFrames.}
    \dayitem{Tuesday}{Glue Studio (Visual ETL) and AWS Glue DataBrew (Data preparation).}
    \dayitem{Wednesday}{Bookmarks, Job tuning (DPUs), and Flex execution.}
    \dayitem{Thursday}{\project{Build a visual ETL pipeline in Glue Studio that joins two tables, drops null values, transforms a timestamp, and loads the clean data into Redshift.}}
    \dayitem{Friday}{\usecase{Evaluate AWS Glue vs. Amazon EMR for a team of business analysts who need to build daily reporting pipelines without managing infrastructure.}}
\end{itemize}

\subsection*{Week 11: Change Data Capture (AWS DMS \& EventBridge)}
\begin{itemize}[label=-]
    \dayitem{Monday}{AWS DMS architecture for continuous replication.}
    \dayitem{Tuesday}{Event-driven data pipelines using Amazon EventBridge.}
    \dayitem{Wednesday}{AWS Lambda for lightweight serverless data transformations.}
    \dayitem{Thursday}{\project{Configure a DMS task with CDC enabled to stream ongoing database changes from an RDS PostgreSQL instance to an S3 Data Lake bucket.}}
    \dayitem{Friday}{\usecase{Architect a near-real-time synchronization pipeline from an on-premise Oracle database to Amazon Redshift for live BI reporting.}}
\end{itemize}

\subsection*{Week 12: Pipeline Orchestration (Amazon MWAA)}
\begin{itemize}[label=-]
    \dayitem{Monday}{Apache Airflow concepts: DAGs, Operators, Tasks, and XComs.}
    \dayitem{Tuesday}{Amazon MWAA (Managed Workflows for Apache Airflow) architecture.}
    \dayitem{Wednesday}{AWS Step Functions for native AWS serverless orchestration.}
    \dayitem{Thursday}{\project{Write a Python DAG for MWAA that: 1. Waits for a file in S3. 2. Triggers an AWS Glue Job. 3. Runs an Athena query. 4. Sends an SNS alert on failure.}}
    \dayitem{Friday}{\usecase{Design a resilient, self-healing orchestration layer for a daily banking reconciliation process spanning 15 dependent systems using MWAA and Step Functions.}}
\end{itemize}

\capstone{Milestone 3}{Automated Batch ELT Pipeline. Use EventBridge to trigger an AWS Step Function state machine upon file upload to S3. The state machine triggers a Glue DataBrew data quality check, runs an EMR Serverless Spark job to calculate complex ML features, and loads the output into Redshift.}

\newpage
% ==========================================
% MILESTONE 4
% ==========================================
\section{Milestone 4: Streaming \& Real-Time Analytics (Month 4)}
\textit{Objective: Conquer unbounded data, Kinesis, Kafka, and real-time visualization.}

\subsection*{Week 13: Message Queuing \& Streaming Ingestion}
\begin{itemize}[label=-]
    \dayitem{Monday}{Decoupling with Amazon SQS and SNS (Fan-out pattern).}
    \dayitem{Tuesday}{Amazon Kinesis Data Streams (KDS): Shards, Partition Keys, and scaling.}
    \dayitem{Wednesday}{Amazon MSK (Managed Streaming for Apache Kafka) vs. Kinesis.}
    \dayitem{Thursday}{\project{Create a Kinesis Data Stream. Write a Python producer script (Boto3) to stream mock IoT temperature sensor data, and a consumer to print the data.}}
    \dayitem{Friday}{\usecase{Architect a global telemetry ingestion point for 5 million mobile games, ensuring exactly-once processing semantics and strict message ordering.}}
\end{itemize}

\subsection*{Week 14: Real-Time Delivery (Kinesis Firehose)}
\begin{itemize}[label=-]
    \dayitem{Monday}{Amazon Kinesis Data Firehose architecture and destinations.}
    \dayitem{Tuesday}{Inline data transformation using AWS Lambda in Firehose.}
    \dayitem{Wednesday}{Format conversion (JSON to Parquet/ORC) on the fly.}
    \dayitem{Thursday}{\project{Connect a Kinesis Firehose delivery stream to your KDS. Configure a Lambda function to uppercase all incoming strings, convert the data to Parquet, and output to S3.}}
    \dayitem{Friday}{\usecase{Design a pipeline to anonymize incoming healthcare HL7 messages in real-time before streaming them into a Redshift cluster for analysis.}}
\end{itemize}

\subsection*{Week 15: Stream Processing (Apache Flink on AWS)}
\begin{itemize}[label=-]
    \dayitem{Monday}{Managed Service for Apache Flink (formerly Kinesis Data Analytics).}
    \dayitem{Tuesday}{Event Time vs. Processing Time and Watermarks in Flink.}
    \dayitem{Wednesday}{Windowing strategies: Tumbling, Sliding, and Session windows.}
    \dayitem{Thursday}{\project{Write a Flink SQL application that reads from a Kinesis Stream, groups user clicks into 5-minute Tumbling Windows, counts the clicks, and writes the output to a new stream.}}
    \dayitem{Friday}{\usecase{Troubleshoot a scenario where a streaming dashboard is showing inaccurate daily aggregations due to mobile devices going offline and uploading data late.}}
\end{itemize}

\subsection*{Week 16: Business Intelligence \& Amazon QuickSight}
\begin{itemize}[label=-]
    \dayitem{Monday}{Amazon QuickSight Architecture and SPICE engine.}
    \dayitem{Tuesday}{Connecting QuickSight to Redshift, Athena, and S3.}
    \dayitem{Wednesday}{Row-level security (RLS) and ML Insights in QuickSight.}
    \dayitem{Thursday}{\project{Connect Amazon QuickSight to your Redshift cluster. Create a real-time dashboard monitoring sales, apply a filter, and set up an ML Anomaly Detection alert.}}
    \dayitem{Friday}{\usecase{Design the BI access layer for an enterprise where 500 analysts need to query centralized data, ensuring regional sales managers can only see data for their specific regions (RLS).}}
\end{itemize}

\capstone{Milestone 4}{Real-Time Fraud Detection. Build a Python script simulating credit card swipes sent to Kinesis Data Streams. Create a Managed Flink application that uses a Sliding Window to detect if a card is used in two different cities within 10 minutes, flags it, routes it via Firehose to Redshift, and displays the live fraud alerts in QuickSight.}


% ==========================================
% MILESTONE 5
% ==========================================
\section{Milestone 5: Machine Learning Operationalization (Month 5)}
\textit{Objective: Embed AI into data pipelines and manage the MLOps lifecycle with SageMaker.}

\subsection*{Week 17: In-Database Machine Learning (Redshift ML)}
\begin{itemize}[label=-]
    \dayitem{Monday}{Redshift ML architecture (integration with SageMaker Autopilot).}
    \dayitem{Tuesday}{Creating models: Regression, Classification, and K-Means.}
    \dayitem{Wednesday}{Model evaluation and scoring directly via SQL queries.}
    \dayitem{Thursday}{\project{Use Redshift ML to train a logistic regression model on sample customer data to predict churn. Run a SQL query using the new model function to predict churn on live data.}}
    \dayitem{Friday}{\usecase{Architect a system to automatically forecast daily inventory needs for 1,000 retail stores using purely SQL-based pipelines to minimize data movement.}}
\end{itemize}

\subsection*{Week 18: Amazon SageMaker Foundations}
\begin{itemize}[label=-]
    \dayitem{Monday}{SageMaker Studio, Notebook Instances, and SageMaker Data Wrangler.}
    \dayitem{Tuesday}{Built-in Algorithms (XGBoost, BlazingText) vs. Custom Containers.}
    \dayitem{Wednesday}{SageMaker Endpoints: Real-time, Serverless, and Batch Transform.}
    \dayitem{Thursday}{\project{Use SageMaker Data Wrangler to clean an S3 dataset. Train an XGBoost model using SageMaker Autopilot. Deploy the model to a serverless endpoint and send a test payload.}}
    \dayitem{Friday}{\usecase{Design a workflow for data scientists to train custom PyTorch models on GPU instances while securely accessing training data stored in a governed S3 data lake.}}
\end{itemize}

\subsection*{Week 19: MLOps \& SageMaker Pipelines}
\begin{itemize}[label=-]
    \dayitem{Monday}{SageMaker Feature Store: Online vs. Offline serving.}
    \dayitem{Tuesday}{SageMaker Pipelines for MLOps orchestration.}
    \dayitem{Wednesday}{SageMaker Model Monitor (Data Drift and Concept Drift).}
    \dayitem{Thursday}{\project{Build a simple SageMaker Pipeline (Preprocess $\rightarrow$ Train $\rightarrow$ Evaluate $\rightarrow$ Register). Execute the pipeline from a Jupyter notebook.}}
    \dayitem{Friday}{\usecase{Architect an MLOps lifecycle that detects data drift in a production recommendation engine and automatically triggers a retraining pipeline using EventBridge.}}
\end{itemize}

\subsection*{Week 20: Unstructured Data \& AI Services}
\begin{itemize}[label=-]
    \dayitem{Monday}{Amazon Rekognition (Images/Video) and Amazon Textract (OCR).}
    \dayitem{Tuesday}{Amazon Comprehend (NLP) and Amazon Translate.}
    \dayitem{Wednesday}{Integrating AWS AI APIs within Lambda and Glue pipelines.}
    \dayitem{Thursday}{\project{Write a Lambda function triggered by an S3 upload that sends a scanned PDF document to Amazon Textract, extracts the key-value pairs, and saves them as JSON back to S3.}}
    \dayitem{Friday}{\usecase{Design a scalable pipeline to process terabytes of uploaded user images, extract text (OCR), tag objects, and make the metadata searchable using OpenSearch.}}
\end{itemize}

\capstone{Milestone 5}{End-to-End MLOps Pipeline. Use MWAA (Airflow) to orchestrate a pipeline that extracts features via Athena, pushes them to SageMaker Feature Store, trains a custom model using SageMaker Training Jobs, registers the model in the Model Registry, and sets up SageMaker Model Monitor for serving skew.}

\newpage
% ==========================================
% MILESTONE 6
% ==========================================
\section{Milestone 6: Security, Governance \& Exam Prep (Month 6)}
\textit{Objective: Secure the perimeter, implement Data Mesh, and prepare for certification.}

\subsection*{Week 21: Security, IAM \& Networking}
\begin{itemize}[label=-]
    \dayitem{Monday}{AWS IAM Deep Dive: Roles, Policies, Boundaries, and SCPs.}
    \dayitem{Tuesday}{VPC Networking: Endpoints (Gateway vs. Interface) and PrivateLink.}
    \dayitem{Wednesday}{AWS KMS (Key Management Service) and AWS Secrets Manager.}
    \dayitem{Thursday}{\project{Create a VPC Endpoint for S3 to ensure traffic doesn't traverse the public internet. Attempt to access the S3 bucket from an EC2 instance without an Internet Gateway.}}
    \dayitem{Friday}{\usecase{Design a highly secure data perimeter for a healthcare provider ensuring no data exfiltration can occur to outside AWS accounts even if IAM credentials are compromised.}}
\end{itemize}

\subsection*{Week 22: Data Governance \& Quality}
\begin{itemize}[label=-]
    \dayitem{Monday}{AWS Lake Formation: Cross-account data sharing and Tag-based access control.}
    \dayitem{Tuesday}{AWS Glue Data Quality: Automated rulesets and alerts.}
    \dayitem{Wednesday}{Data Lineage with Amazon SageMaker ML Lineage Tracking.}
    \dayitem{Thursday}{\project{Implement an AWS Glue Data Quality ruleset within an ETL job to fail the pipeline if a required column is more than 5\% null. Route failure alerts to SNS.}}
    \dayitem{Friday}{\usecase{Architect a Data Mesh for a global conglomerate with 5 independent domains, ensuring centralized governance via Lake Formation but decentralized data engineering.}}
\end{itemize}

\subsection*{Week 23: Compliance \& Data Privacy (Amazon Macie)}
\begin{itemize}[label=-]
    \dayitem{Monday}{Amazon Macie: Automated PII discovery and classification in S3.}
    \dayitem{Tuesday}{Masking, Tokenization, and Crypto-based hashing in pipelines.}
    \dayitem{Wednesday}{AWS CloudTrail and AWS Config for auditing data access.}
    \dayitem{Thursday}{\project{Enable Amazon Macie on an S3 bucket. Upload a CSV containing mock credit card numbers. Review the Macie findings report and trigger a Lambda function to quarantine the file.}}
    \dayitem{Friday}{\usecase{Design a pipeline that ingests third-party customer data, automatically identifies PII using Macie, tokenizes it using a KMS key in Glue, and stores the safe version in the data warehouse.}}
\end{itemize}

\subsection*{Week 24: Final Certification Preparation}
\begin{itemize}[label=-]
    \dayitem{Monday}{Review AWS Certified Data Engineer – Associate/Professional Domains.}
    \dayitem{Tuesday}{Whitepaper Review: Big Data Analytics Options on AWS.}
    \dayitem{Wednesday}{Full-length Mock Exam 1 (Identify weak domains).}
    \dayitem{Thursday}{Targeted review of weak domains (e.g., Kinesis vs. MSK, Redshift optimization).}
    \dayitem{Friday}{Full-length Mock Exam 2.}
\end{itemize}

\capstone{Milestone 6}{Secure, Governed Data Mesh. Build a multi-account architecture using AWS Organizations. Account A ingests streaming PII data and uses Macie/Lambda to tokenize it. Account B (Analytics) hosts Redshift. Account C (Governance) uses Lake Formation to enforce cross-account tag-based security and VPC Endpoints to lock down all traffic.}

\newpage
% ==========================================
% TRACKING TIMETABLE
% ==========================================
\section{6-Month Progress Tracking Timetable}

\begin{center}
\renewcommand{\arraystretch}{1.5}
\begin{longtable}{@{} c l p{7cm} c @{}}
\toprule
\textbf{Week} & \textbf{Primary Focus} & \textbf{Key Project / Capstone} & \textbf{Status} \\ 
\midrule
\endfirsthead

\toprule
\textbf{Week} & \textbf{Primary Focus} & \textbf{Key Project / Capstone} & \textbf{Status} \\ 
\midrule
\endhead

\bottomrule
\endfoot

\bottomrule
\endlastfoot

1 & Amazon S3 \& Migration & S3 Lifecycle \& DataSync & $\square$ \\
2 & Amazon RDS \& Aurora & Aurora Serverless Failover Simulation & $\square$ \\
3 & Amazon DynamoDB & DynamoDB Single-Table Design \& GSI & $\square$ \\
4 & Hybrid Ingestion \& Caches & \textbf{Cap 1: E-Commerce Multi-Tier Data} & $\square$ \\
\midrule
5 & Redshift Architecture & COPY from S3 to Redshift Serverless & $\square$ \\
6 & Redshift Optimization & Dist/Sort Keys Benchmark & $\square$ \\
7 & Athena \& Lakehouse & Glue Crawler \& Athena vs Spectrum & $\square$ \\
8 & AWS Lake Formation & \textbf{Cap 2: Enterprise Data Lakehouse} & $\square$ \\
\midrule
9 & Amazon EMR & Spark Job on EMR Cluster & $\square$ \\
10 & AWS Glue ETL & Glue Studio Visual Pipeline & $\square$ \\
11 & AWS DMS \& EventBridge & RDS to S3 CDC via DMS & $\square$ \\
12 & Amazon MWAA (Airflow) & \textbf{Cap 3: Automated Batch ELT} & $\square$ \\
\midrule
13 & Amazon Kinesis Data Streams & IoT Producer/Consumer Scripting & $\square$ \\
14 & Kinesis Data Firehose & Firehose inline Lambda format conversion & $\square$ \\
15 & Managed Apache Flink & Flink Tumbling Windows on Streams & $\square$ \\
16 & Amazon QuickSight & \textbf{Cap 4: Real-Time Fraud Telemetry} & $\square$ \\
\midrule
17 & Redshift ML & Churn Prediction SQL Models & $\square$ \\
18 & SageMaker Foundations & Autopilot Deployment \& Endpoint Test & $\square$ \\
19 & SageMaker Pipelines & SageMaker Pipeline execution & $\square$ \\
20 & AI APIs (Textract/Rekognition)& \textbf{Cap 5: End-to-End MLOps Pipeline} & $\square$ \\
\midrule
21 & IAM \& VPC Endpoints & S3 Gateway Endpoint implementation & $\square$ \\
22 & Glue Data Quality & Automated DQ Rulesets \& Alerts & $\square$ \\
23 & Amazon Macie & PII Discovery \& Lambda Quarantine & $\square$ \\
24 & Exam Prep & \textbf{Cap 6: Secure Data Mesh Design} & $\square$ \\

\end{longtable}
\end{center}

\vspace{1cm}
\begin{center}
    \textit{"Consistency is the architecture of mastery."} \\
    Best of luck on your 6-month journey to becoming a Master AWS Data Engineer!
\end{center}

\end{document}